% Options for packages loaded elsewhere
\PassOptionsToPackage{unicode}{hyperref}
\PassOptionsToPackage{hyphens}{url}
\documentclass[
  a4paper,
]{article}
\usepackage{xcolor}
\usepackage[margin=25mm]{geometry}
\usepackage{amsmath,amssymb}
\setcounter{secnumdepth}{-\maxdimen} % remove section numbering
\usepackage{iftex}
\ifPDFTeX
  \usepackage[T1]{fontenc}
  \usepackage[utf8]{inputenc}
  \usepackage{textcomp} % provide euro and other symbols
\else % if luatex or xetex
  \usepackage{unicode-math} % this also loads fontspec
  \defaultfontfeatures{Scale=MatchLowercase}
  \defaultfontfeatures[\rmfamily]{Ligatures=TeX,Scale=1}
\fi
\usepackage{lmodern}
\ifPDFTeX\else
  % xetex/luatex font selection
\fi
% Use upquote if available, for straight quotes in verbatim environments
\IfFileExists{upquote.sty}{\usepackage{upquote}}{}
\IfFileExists{microtype.sty}{% use microtype if available
  \usepackage[]{microtype}
  \UseMicrotypeSet[protrusion]{basicmath} % disable protrusion for tt fonts
}{}
\makeatletter
\@ifundefined{KOMAClassName}{% if non-KOMA class
  \IfFileExists{parskip.sty}{%
    \usepackage{parskip}
  }{% else
    \setlength{\parindent}{0pt}
    \setlength{\parskip}{6pt plus 2pt minus 1pt}}
}{% if KOMA class
  \KOMAoptions{parskip=half}}
\makeatother
\usepackage{longtable,booktabs,array}
\newcounter{none} % for unnumbered tables
\usepackage{calc} % for calculating minipage widths
% Correct order of tables after \paragraph or \subparagraph
\usepackage{etoolbox}
\makeatletter
\patchcmd\longtable{\par}{\if@noskipsec\mbox{}\fi\par}{}{}
\makeatother
% Allow footnotes in longtable head/foot
\IfFileExists{footnotehyper.sty}{\usepackage{footnotehyper}}{\usepackage{footnote}}
\makesavenoteenv{longtable}
\usepackage{graphicx}
\makeatletter
\newsavebox\pandoc@box
\newcommand*\pandocbounded[1]{% scales image to fit in text height/width
  \sbox\pandoc@box{#1}%
  \Gscale@div\@tempa{\textheight}{\dimexpr\ht\pandoc@box+\dp\pandoc@box\relax}%
  \Gscale@div\@tempb{\linewidth}{\wd\pandoc@box}%
  \ifdim\@tempb\p@<\@tempa\p@\let\@tempa\@tempb\fi% select the smaller of both
  \ifdim\@tempa\p@<\p@\scalebox{\@tempa}{\usebox\pandoc@box}%
  \else\usebox{\pandoc@box}%
  \fi%
}
% Set default figure placement to htbp
\def\fps@figure{htbp}
\makeatother
\setlength{\emergencystretch}{3em} % prevent overfull lines
\providecommand{\tightlist}{%
  \setlength{\itemsep}{0pt}\setlength{\parskip}{0pt}}
\usepackage{bookmark}
\IfFileExists{xurl.sty}{\usepackage{xurl}}{} % add URL line breaks if available
\urlstyle{same}
\hypersetup{
  hidelinks,
  pdfcreator={LaTeX via pandoc}}

\author{}
\date{}

\begin{document}

\section{How Inflation Ends Is How Matter Generation Begins: An N-Body
Embedding of the Universal Interaction with No Timing Rule, and an
End-to-End Genesis
Run}\label{how-inflation-ends-is-how-matter-generation-begins-an-n-body-embedding-of-the-universal-interaction-with-no-timing-rule-and-an-end-to-end-genesis-run}

\textbf{Author:} Noriaki Kihara (WF System Co., Ltd.) \textbf{Date:}
2026-08-06 \textbf{Version:} v1 \textbf{Version DOI:}
\href{https://doi.org/10.5281/zenodo.21809815}{10.5281/zenodo.21809815}
\textbf{Concept DOI:}
\href{https://doi.org/10.5281/zenodo.21809814}{10.5281/zenodo.21809814}

\begin{center}\rule{0.5\linewidth}{0.5pt}\end{center}

\subsection{Abstract}\label{abstract}

Two preceding papers each supplied half of the story. The onset-mode
paper {[}2{]} gave the way a universe begins (only unstable
self-consistent closures undergo latent exponential amplification,
ending in a three-direction metastable structure) --- but there the
dynamics stops and there is no matter. The two-body paper {[}1{]} gave
the generation mechanism (a universal inelastic vertex whose only input
is the waveforms) --- but its collisions were loop counters, with when
and where unanswered. This paper closes the gap \textbf{without
introducing any timing rule}. Embedding the two-body vertex into the
N-body relational wave system (the \(M=N(N-1)/2\) relational waves of
the complete graph \(\times\) a multi-slice register axis) with the
interaction always on, we measure the following. (1) The apparent epoch
of generation is set by a \textbf{two-factor law} --- generation rate =
condensation gate \(\times\) \(Cf^2\) (gate ratio 2.74 measured). And
the condensation is precisely the endpoint of the onset paper:
\textbf{that the ending of the previous paper is the opening condition
of this one is a measurement, not an assumption.} (2) The two-body laws
hold unchanged in the many-body system: the ignition law transplants
with power \(p=2.0011\) and coefficient constant to 0.4\%, and the
pair-structure census holds up to the hair (winding-number) complete
version (sum-rule exclusivity ratio \(10^{289}\), hair lock 0.9997). (3)
Running vacuum \(\to\) latent exponential amplification \(\to\)
three-direction condensation \(\to\) matter precipitation under a single
dynamics with no switches (GENESIS), we discover \textbf{seed
unification}: the perturbation that ignites inflation and the seed that
becomes matter are one and the same (\(\varepsilon\)-law agreement to
2\%). (4) Locality does not live on the graph --- concentration onto the
complete graph leaves the generation rate unchanged (ratios 1.28, 1.02)
--- but on the register axis (response \(\times\)\textasciitilde19.6).
Products are pinned at their birth position, and harmonic-closure
dispersion provides shape-invariant translation (1D deviation below
\(3\times10^{-4}\) cells; 3D deviation 0.0000). (5) The generation
direction is a phase lottery of the seed (\(P(+)=0.52\) over 50 random
phases), and net matter arises by rectification through the \(Cf^2\)
autocatalysis (80\% of the ensemble grows net). As conditional
derivations we present the mediated-vertex uniqueness theorem (closure
conservation forces \(g_2=-g_1\) \textbf{and} pair symmetry of the
mediated strengths --- the two-body uniqueness is its \(M=2\) special
case) and the identity of the computational structure (the adjacency sum
folds exactly into two edge\(\to\)node\(\to\)edge passes; the operation
count is linear in the state size, the theoretical lower bound, and what
guarantees this is anonymity itself). We report with assumptions,
conditional derivations, and measurements strictly separated, and we
measure the value of the assumption by the demonstrated connection ---
that the physics of two independently published papers reproduces
simultaneously as different chapters of a single dynamics with no
switches.

\begin{center}\rule{0.5\linewidth}{0.5pt}\end{center}

\subsection{1. The Character of This Paper --- Separating Assumption,
Derivation, and
Measurement}\label{the-character-of-this-paper-separating-assumption-derivation-and-measurement}

This paper neither proposes a new interaction nor proves one correct. It
takes the universal inelastic vertex --- posited as an explicit
assumption and conditionally uniquified in the two-body paper {[}1{]}
--- embeds it into the N-body relational wave system with no timing rule
whatsoever, and measures what happens.

\begin{itemize}
\tightlist
\item
  \textbf{(a) Assumptions}: the inheritance of the vertex from the
  two-body paper, and the design choices of the embedding (Section 4).
\item
  \textbf{(b) Conditional derivations}: the mediated-vertex uniqueness
  theorem and the identity of the computational structure (Section 5).
\item
  \textbf{(c) Measurements}: the resolution of timing, the
  transplantation of the two-body laws, the end-to-end genesis run,
  position and motion, and the direction lottery with rectification
  (Sections 6--10).
\end{itemize}

What the reader can verify are (b) and (c); (a) is the starting point.

\subsection{2. Claims}\label{claims}

\textbf{First --- the motivation and design principle of this paper ---
we do not choose by hand when the interaction happens.} The two
preceding papers each gave half. The onset paper gave the way a universe
begins, but the dynamics stops at its endpoint (the three-direction
metastable structure) with no matter. The two-body paper gave the
generation mechanism, but its collisions were loop counters. Filling
this gap with a rule such as ``invoke the interaction after
condensation'' would be an IF statement, and the anonymity the two-body
paper preserved would collapse in the many-body setting. The design of
this paper is single: \textbf{the interaction acts at all times}.
Eventness is an output, not an input.

\textbf{Second, we enumerate all assumptions (Section 4).} (i) The
linear dynamics is inherited unmodified from Papers 7 and 8 (read-only
import, SHA recorded). (ii) The vertex extends the uniquified two-body
form to an adjacency sum --- each relational wave couples through the
simple sum of two-wave vertices with the relational waves sharing a body
(no vertex of three or more waves). (iii) The state is a multi-slice
register \(C\in\mathbb{C}^{M\times N_{reg}}\) with a shared readout of
the combined spectrum. (iv) No new constants --- the strength reuses the
readout \(R\). \textbf{The four items are design choices, not
derivations.}

\textbf{Third, the mediated-vertex uniqueness theorem (conditional
derivation, Section 5.1).} Within the assumed class, closure
conservation forces \(g_2=-g_1\) \textbf{and} pair symmetry of the
mediated strengths. The two-body \(g_2=-g_1\) is the \(M=2\) special
case of this theorem. What closure demands of the strengths is pair
symmetry only; the arithmetic-mean rule is an additional experimental
fact not derived from closure. Extension theorem T-L:
register-point-dependent pair-symmetric strengths \(R_{ee'}(n)\) also
conserve closure exactly --- locally varying interaction strength is
compatible with the axiom (machine verified).

\textbf{Fourth, the identity of the computational structure (conditional
derivation, Section 5.2).} The adjacency sum of the vertex folds exactly
into a two-pass edge\(\to\)node\(\to\)edge aggregation (vertex moments
and self-subtraction). The operation count is linear in the state size,
\(O(M\cdot N_{reg}\log N_{reg})\) --- no all-pairs checks \(O(N^4)\), no
adjacency scans \(O(N^3)\) --- reaching the theoretical lower bound.
What makes this folding possible is anonymity itself (one law for all
pairs, no pair-specific constants).

\textbf{Fifth, Measurement 1: the timing problem resolves into the
two-factor law (Section 6).} Generation rate = condensation gate
\(\times\) \(Cf^2\). The measured gate ratio is 2.74. Condensation is
the endpoint of the onset paper --- the three-direction metastable
structure. That the previous paper's ending is this paper's opening
condition is a measurement, not an assumption.

\textbf{Sixth, Measurement 2: the two-body laws hold unchanged in the
many-body system (Section 7).} The ignition law reproduces with power
\(p=2.0011\) and constant coefficient (0.4\%). The census holds up to
the hair-complete version: sum-rule exclusivity (\(10^{-289}\) for
unpredicted bins), hair lock 0.9997, simultaneous rise from exact zero.

\textbf{Seventh, Measurement 3: the end-to-end genesis run GENESIS
(Section 8).} Vacuum \(\to\) latent exponential amplification \(\to\)
three-direction condensation \(\to\) matter precipitation under one
dynamics with no switches. Discovery: \textbf{seed unification} --- the
perturbation that ignites inflation and the seed that becomes matter are
one and the same.

\textbf{Eighth, Measurement 4: position and motion (Section 9).}
Locality does not live on the graph (concentration ineffective; ratios
1.28, 1.02). Locality lives on the register axis (response
\(\times\)\textasciitilde19.6). Products are pinned at their birth
position, and motion is provided by the harmonic-closure dispersion
\(\omega_k=k\,\omega_1\) {[}6{]} (shape-invariant translation
demonstrated in 1D/3D).

\textbf{Ninth, Measurement 5: the direction lottery and rectification
(Section 10).} The generation direction is set by the seed's phases
(\(P(+)=0.52\) over 50 random phases). The instantaneous sign predictor
scores 0.70 --- above chance (0.5) but short of the preregistered
criterion 0.9; the determinism of the lottery is only partially
confirmed (a complete instantaneous predictor is an open problem). Net
matter arises not from a hand-placed symmetry breaking but from
rectification by the \(Cf^2\) autocatalysis; 80\% of the ensemble grew
net. This is a measurement of a mechanism inside the model; we do not
claim it explains the matter excess of the real universe.

\textbf{Tenth, we record a null result (Section 11).} The selectivity of
the E8-related special values (\(N_{reg}=124,137,144,248\), etc.) is
null on every instrument tested.

\textbf{Eleventh, we record the full process of refutation and
correction (Section 12)}, including three families of instrument error
(conservation criteria, generation-rate predictors, the motion centroid
gauge).

\textbf{Finally, what this paper does not claim, and the falsification
conditions (Section 15).}

\subsection{3. Motivation --- the Two Papers Were Not
Connected}\label{motivation-the-two-papers-were-not-connected}

The conclusion of the onset paper {[}2{]} is this: only an unstable
self-consistent closure (a single coherent condensate) undergoes latent
exponential amplification, and nothing amplifies out of white chaos. The
endpoint of amplification is the three-direction metastable structure,
where the linear dynamics completes. \textbf{Spacetime was made. But
there is no matter and nothing happens.}

The conclusion of the two-body paper {[}1{]} is this: a single IF-free
map seeing only waveforms generates fermionic structure --- induced,
autocatalytic, pair-correlated. \textbf{We know how to make matter. But
there is no universe.} The two-body ``collision'' was a loop counter;
when, where, and why collisions occur lay outside the dynamics.

This gap was not an incompleteness but a candidate fatal wound. Choosing
the timing by hand --- ``invoke the interaction after inflation'' --- is
an IF statement of the same rank as discriminating particle species, and
the anonymity the two-body paper defended would collapse at the moment
of embedding. A model that invokes its interaction only when convenient
has no standing to speak about generation. Hence there is only one
design: embed with the interaction always on, and measure whether
``when'' is output as physics. That is this paper.

\subsection{4. Model and Assumptions}\label{model-and-assumptions}

\subsubsection{4.1 The relational wave graph
(inherited)}\label{the-relational-wave-graph-inherited}

The fundamental variables do not live on bodies (particles). Each pair
of the \(N\) bodies (nodes) carries one relational wave
\(z_e\in\mathbb{C}\) (an edge; \(M=N(N-1)/2\) in total). Nodes
themselves hold no variables; they are aggregation points for ledgers.
Two relational waves are eligible to interact only when they share a
body (adjacency) --- the shared body mediates (Figure 1). The linear
dynamics (generator \(K_{ef}=\sin(\theta_f-\theta_e)\), Cayley step;
self-consistent closure, latency, burst, three-direction generation) is
inherited unmodified from the original code of Papers 7 and 8 {[}3,4{]}
(SHA-256 recorded, read-only import).

\begin{figure}
\centering
\pandocbounded{\includegraphics[keepaspectratio,alt={Figure 1 relational wave graph}]{figs/fig_graph_N6_v1.png}}
\caption{Figure 1 relational wave graph}
\end{figure}

\textbf{Figure 1}: the relational wave graph at N=6. (a) Bodies = nodes;
relational waves = edges (M=15 dynamical variables). (b) Adjacency =
sharing a body: the interaction partners of z12 (red) are only the 8
edges sharing body 1 or body 2 (blue); the 6 edges sharing no body (gray
dashed) do not couple directly. (c) The node ledger: the sum \(A_1\) of
the relational waves meeting body 1.

\subsubsection{4.2 Assumptions of the
embedding}\label{assumptions-of-the-embedding}

{\def\LTcaptype{none} % do not increment counter
\begin{longtable}[]{@{}
  >{\raggedright\arraybackslash}p{(\linewidth - 4\tabcolsep) * \real{0.2000}}
  >{\raggedright\arraybackslash}p{(\linewidth - 4\tabcolsep) * \real{0.4000}}
  >{\raggedright\arraybackslash}p{(\linewidth - 4\tabcolsep) * \real{0.4000}}@{}}
\toprule\noalign{}
\begin{minipage}[b]{\linewidth}\raggedright
\#
\end{minipage} & \begin{minipage}[b]{\linewidth}\raggedright
Assumption
\end{minipage} & \begin{minipage}[b]{\linewidth}\raggedright
Content
\end{minipage} \\
\midrule\noalign{}
\endhead
\bottomrule\noalign{}
\endlastfoot
(i) & Inherited linear dynamics & Read-only import of the Papers 7/8
originals. Acts independently on each slice (register harmonic) \\
(ii) & Vertex = adjacency sum of two-wave vertices & Each edge \(e\)
couples through the simple sum over adjacent edges:
\(\delta z_e = i\sum_{e'\sim e} R_{ee'}(|z_{e'}|^2 z_e - z_{e'}^2\bar z_e)\).
No vertex of three or more waves \\
(iii) & Multi-slice register & State
\(C\in\mathbb{C}^{M\times N_{reg}}\). The vertex acts pointwise in
register position space. The readout is shared over the combined
spectrum: \(\theta=\operatorname{atan2}(\sqrt{P_f},\sqrt{P_b})\),
\(R=\sin^2\theta\) \\
(iv) & No new constants & The strength reuses \(R\). No discriminating
branches (the implementation's \texttt{r\textgreater{}0} is a
zero-strength numerical skip) \\
\end{longtable}
}

The four items are design choices. (ii) two-wave and (iv) anonymity are
direct inheritances of the two-body assumptions; (i) and (iii) are
embedding choices specific to this paper (Figure 2).

\begin{figure}
\centering
\pandocbounded{\includegraphics[keepaspectratio,alt={Figure 2 architecture}]{figs/fig_mb1_architecture_v1.png}}
\caption{Figure 2 architecture}
\end{figure}

\textbf{Figure 2}: the state is the product lattice of the graph axis (M
relational waves) \(\times\) the register axis (\(N_{reg}\) points). The
linear part acts per slice (column); the vertex acts per register point
through the node ledgers.

\subsection{5. Conditional Derivations}\label{conditional-derivations}

\subsubsection{5.1 The mediated-vertex uniqueness
theorem}\label{the-mediated-vertex-uniqueness-theorem}

For the most general form in the assumed class (two-wave, pointwise,
cubic, common-phase invariant, no self-scattering),
\(\delta z_e = i\sum_{e'\sim e}(g_1 R_{ee'}|z_{e'}|^2 z_e + g_2 R'_{ee'} z_{e'}^2\bar z_e)\),
conservation of the zero-square-sum closure \(\sum_e z_e^2\) (per
register point) forces:

\[g_2 = -g_1 \quad\text{and}\quad R_{ee'} = R_{e'e}\ (\text{pair symmetry}).\]

The proof is by pairwise cancellation of the first variation: the
contribution of the edge pair \((e,e')\) is
\(g_1 R_{ee'}|z_{e'}|^2 z_e^2 + g_2 R'_{ee'} z_{e'}^2 |z_e|^2\) plus its
\(e\leftrightarrow e'\) exchange, and it vanishes identically only when
the coefficients are opposite and the strengths pair-symmetric. Three
points. First, this is uniqueness within the class. Second, the two-body
\(g_2=-g_1\) {[}1{]} is the \(M=2\) special case (the adjacency sum has
one partner). Third, closure demands of the strengths \textbf{pair
symmetry only}: as the mixing rule when two edges have different \(R\),
both the arithmetic mean and the product conserve closure exactly
(machine test T5). That the actual dynamics is consistent with the
arithmetic mean is an additional experimental fact.

\textbf{Extension theorem T-L}: since pair symmetry holds per point
\(n\), register-point-dependent pair-symmetric strengths \(R_{ee'}(n)\)
also conserve closure exactly. Machine check: drift
\(1.0\times10^{-17}\) (machine zero) for random \(n\)-dependent
pair-symmetric \(R\); \(1.7\times10^{-2}\) (broken) when asymmetric.
\textbf{Locally varying interaction strength is compatible with the
axiom} --- the axiom does not forbid locality.

\subsubsection{5.2 The identity of the computational structure ---
anonymity guarantees
computability}\label{the-identity-of-the-computational-structure-anonymity-guarantees-computability}

Naively, evaluating the adjacency sum scans \(2(N-2)\) adjacent edges
per edge, costing \(O(M\cdot N)=O(N^3)\) per register point. But by (ii)
two-wave additivity (no cross terms between partners) and the
separability of the pair-symmetric mean strength, the adjacency sum
folds \textbf{exactly}: for edge \(e=(u,v)\),

\[\sum_{e'\sim e} |z_{e'}|^2 = A_u + A_v - 2|z_e|^2, \qquad A_v \equiv \sum_{e\ni v}|z_e|^2\]

(likewise for the four ledgers including the \(z^2\)-type and
\(R\)-weighted ones). The computation is two passes --- ``edges \(\to\)
node aggregation'' and ``nodes \(\to\) edge read-back with
self-subtraction'' --- each \(O(M)\). This is an identity, not a
mean-field approximation; agreement with the naive implementation is
verified to machine precision. The total per step is
\(O(M\cdot N_{reg}\log N_{reg})\) --- the same order as the trivial
lower bound in the RAM model of having to read the whole state once per
step, and in this sense the theoretical lower bound.

The physical meaning: information flows edge \(\to\) shared node \(\to\)
edge. \textbf{The interaction is not a pairwise direct check but
aggregation and redistribution mediated by the bodies} --- the \(O(M)\)
algorithm is the literal implementation of ``mediated vertex.'' And what
makes the folding possible is anonymity itself: give each pair its own
coupling table \(R_{ee'}\) (violating anonymity) and separability dies,
returning the cost to \(O(N^3)\). A physical principle directly
guarantees a computational property.

\textbf{Measured conserved quantities} (under the discipline of deriving
the invariant before judging): relative norm drift
\(3.1\times10^{-14}\); vertex-only pointwise closure drift
\(7.3\times10^{-11}\); linear-only slice closure drift
\(1.4\times10^{-15}\) (each judged only where it is an invariant; see
Section 12). The ON/OFF separation of generated content is 11 orders;
the ON amount is consistent with the value derived from the ignition law
within a factor 4.2.

\subsection{6. Measurement 1: the Two-Factor Law --- Resolving the
Timing
Problem}\label{measurement-1-the-two-factor-law-resolving-the-timing-problem}

With the always-on embedding, ``when generation happens'' did not
vanish; it reappeared as physics:

\[\text{generation rate} = \text{condensation gate} \times C f^2.\]

\textbf{The condensation gate} (Figure 3a): in a control with identical
norm and identical seed, with the pump placed either in the condensed
state (the three-direction metastable structure) or in a non-condensed
state (a non-coherent state of the same energy), the generation rates
are \(6.30\times10^{-8}\) vs.~\(2.30\times10^{-8}\) --- \textbf{gate
ratio 2.74}. Condensation is precisely the endpoint of the onset paper
{[}2{]}.

\textbf{Era modulation} (Figure 3b): the windowed generation rate of the
end-to-end run (Section 8) opens monotonically --- latency
\(3.3\times10^{-7}\) \(\to\) burst \(1.4\times10^{-6}\) \(\to\)
post-condensation \(6.4\times10^{-6}\) --- a factor
\textbf{\(\times\)\textasciitilde19.5} from latency to
post-condensation. The interaction acts at all times, yet the apparent
epoch of generation is made by the age of the universe.

\begin{figure}
\centering
\pandocbounded{\includegraphics[keepaspectratio,alt={Figure 3 two-factor gate}]{figs/fig_mb2_gate_v1.png}}
\caption{Figure 3 two-factor gate}
\end{figure}

\textbf{Figure 3}: (a) the condensation-gate control. (b) era modulation
of the end-to-end run. No timing rule was written; timing was output.

\textbf{Reading (explicitly an interpretation)}: this measurement
suggests that it is not that ``after inflation ends, matter generation
begins as a separate phenomenon.'' The end of expansion (linear side:
saturation of the unstable growth, settling into the three-direction
structure) and the onset of generation (nonlinear side: opening of the
condensation gate) can be read as \textbf{two observables of the same
condensation transition} --- spacetime formation and matter generation
are not two serial processes but the outer geometric face and the inner
generative face of one transition. The reading is an interpretation; the
two quantities supporting it (gate ratio, era modulation) are
measurements.

\subsection{7. Measurement 2: Transplantation Tests of the Two-Body
Laws}\label{measurement-2-transplantation-tests-of-the-two-body-laws}

\textbf{The ignition law} (Figure 4): over a four-decade sweep of the
seed fraction \(f_0\), the initial generation rate obeys
\(\mathrm{rate}=Cf^p\) with \(p=2.0011\) and a coefficient spread of
0.37\%. Across an engine revision (stage 2 \(\to\) stage 3
shared-readout), the law is unchanged: \(p=2.0015\), \(C=5.88\). The
two-body autocatalytic ignition law {[}1{]} keeps its form inside the
many-body sea.

\begin{figure}
\centering
\pandocbounded{\includegraphics[keepaspectratio,alt={Figure 4 ignition transplant}]{figs/fig_mb3_ignition_v1.png}}
\caption{Figure 4 ignition transplant}
\end{figure}

\textbf{Figure 4}: (a) the many-body ignition law rate=C·f\^{}p
(log-log, four decades). (b) constancy of the coefficient.

\textbf{The pair-structure census, hair-complete} (Figure 5): the
preregistered census with narrow-band pumps and a single seed was run on
the two-dimensional ledger of register (\(k\)) and hair (\(\eta\)
winding \(m\)). The hair-resolved power of the partner band \(k=3\)
appears only at the hair dictated by the sum rule
\(m^* = 2m_B - m_s = 3\) (\textbf{exclusivity ratio
\(6.1\times10^{289}\)} against unpredicted hairs at machine zero), with
pair-correlation lock 0.9997. The partner band rises from exact zero
simultaneously with the sidebands. The ledger is machine-derived from a
single vertex application --- nothing hand-set. The pair-correlation
lock decays slowly over long times (0.988 at the tenth window; Figure
5b) --- the cause is unresolved and recorded as an open problem.

\begin{figure}
\centering
\pandocbounded{\includegraphics[keepaspectratio,alt={Figure 5 hair-complete census}]{figs/fig_mb4_census_v1.png}}
\caption{Figure 5 hair-complete census}
\end{figure}

\textbf{Figure 5}: (a) the hair ledger: only the predicted \(m^*=3\)
generates. (b) persistence and slow decay of the lock (open problem).

\subsection{8. Measurement 3: GENESIS --- the End-to-End Genesis
Run}\label{measurement-3-genesis-the-end-to-end-genesis-run}

From an initial state of vacuum (the unstable branch of the
two-chirality vacuum pair) plus a tiny seed \(\delta\), we ran the
single dynamics for 4000 steps (Figure 6). There are no switches, no era
changes, no parameter changes.

\begin{figure}
\centering
\pandocbounded{\includegraphics[keepaspectratio,alt={Figure 6 GENESIS}]{figs/fig_mb8_genesis_v1.png}}
\caption{Figure 6 GENESIS}
\end{figure}

\textbf{Figure 6}: top: the slice-2 vacuum's unstable branch amplifies
exponentially after latency (crossing \(t=180\)) and condenses into the
three-direction metastable structure. Middle: the one-step residual
\(r\) (clock). Bottom: the moving-window generation rate opens with the
era (\(\times\)\textasciitilde19.5) --- the end-to-end demonstration of
the two-factor law.

\textbf{Discovery 1: seed unification.} The pre-run prediction assumed
inflation would be ignited by the system's intrinsic floor (the parent
eigenmode residual \(2.1\times10^{-13}\) {[}2{]}) and registered the
crossing at 1161--1171. The measurement was 555 (the \(\delta=10^{-6}\)
series) --- far outside the band. Diagnosis: \textbf{the seed \(\delta\)
injected for matter was itself igniting inflation.} Substituting
\(\varepsilon=\delta=10^{-6}\) into the onset paper's
\(\varepsilon\)-law (crossing linear in \(-\ln\varepsilon\) with slope
\(1/\lambda\)) predicts about 543, agreeing with the measured 555 to
2\%. The seed that ignites inflation and the seed that becomes matter
are not separate devices but one and the same perturbation. A failed
prediction turned into the discovery of the connection. Figure 10
summarizes the causal structure up to this point.

\begin{figure}
\centering
\pandocbounded{\includegraphics[keepaspectratio,alt={Figure 10 causal diagram of the genesis run}]{figs/fig_mb9_causal_v1.png}}
\caption{Figure 10 causal diagram of the genesis run}
\end{figure}

\textbf{Figure 10}: the causal diagram of the genesis run. From the
single perturbation \(\delta\), two branches: the geometry branch
(unstable departure \(\to\) latent exponential amplification \(\to\)
three-direction condensation \(\to\) gate opens) and the seed branch
(odd-parity seed \(f_0\) \(\to\) autocatalytic factor \(Cf^2\)). They
merge at generation rate \(=G_{cond}\times Cf^2\), followed by matter
precipitation, localization, motion, and pair creation. There is no
timing rule anywhere.

\textbf{Consequence: a sea with a small seed is effectively inert.} At
\(\delta=10^{-6}\) (\(f_0=10^{-12}\)) generation over 3000 steps stays
below the measurement limit (\(|g|<10^{-11}\)), while at
\(\delta=10^{-2}\) it is clear (\(6.4\times10^{-6}\)). The two-body
consequence --- ignition time scaling as \(1/f_0^2\) --- appears
unchanged in the many-body universe.

\subsection{9. Measurement 4: Position and
Motion}\label{measurement-4-position-and-motion}

\textbf{Locality does not live on the graph} (Figure 7a). Concentrating
the seed onto a star subset of the complete graph (concentration
\(\times\)\textasciitilde6) leaves the generation rate unchanged from a
uniform seed (ratios 1.28, 1.02). The complete graph \(K_N\) has no
distance, and the vertex sees adjacency only through node-ledger
aggregation, so ``closeness'' on the graph does not exist in the
physics. Stated precisely, this is not a permanent theorem but the
measurement that ``in the present complete-graph,
anonymous-adjacency-sum embedding, graph concentration did not change
the generation rate'' --- the claim of this paper is that the primitive
complete graph is not position space, and the possibility remains open
that an induced distance structure appears on the graph after the
unification of the three-direction frame with the register (Section
14-4).

\textbf{Locality lives on the register axis} (Figure 7b). Localizing a
seed of identical power into a packet in register position space
(concentration \(\times\)\textasciitilde8, effective width PR=2) raises
the generation response \textbf{\(\times\)\textasciitilde19.6}. Moreover
the register distribution of the products stays at PR=2, pinned at the
birth position throughout growth --- \textbf{particles have positions}.

\begin{figure}
\centering
\pandocbounded{\includegraphics[keepaspectratio,alt={Figure 7 the address of locality}]{figs/fig_mb5_locality_v1.png}}
\caption{Figure 7 the address of locality}
\end{figure}

\textbf{Figure 7}: (a) graph concentration is ineffective. (b) register
concentration gives \(\times\)\textasciitilde19.6. The address of
locality is the register axis, not the graph.

\textbf{Motion} (Figure 8). The linear part has no dispersion (all
harmonics equal speed), so a pinned packet does not move. Introducing
the harmonic-closure dispersion \(\omega_k = k\,\omega_1\) --- the
dispersion relation the inverse-square-law paper {[}6{]} derived from
the geometry of acceleration --- as a register clock is nothing but a
translation operator, and the packet moves rigidly at constant velocity.
Measured: in 1D all 8 recorded points deviate from prediction by at most
\(3\times10^{-4}\) cells, and at integer-cell shifts the waveform
returns exactly to its initial shape (PR 1.000). In 3D (\(8^3\)
register, three independent velocities) all recorded points deviate by
0.0000, PR constant 3.39. This is rigid translation with the vertex ON
and the seed present.

\begin{figure}
\centering
\pandocbounded{\includegraphics[keepaspectratio,alt={Figure 8 motion}]{figs/fig_mb7_motion_v1.png}}
\caption{Figure 8 motion}
\end{figure}

\textbf{Figure 8}: (a) 1D translation (mod-8 display). (b) 3D uniform
linear motion.

\subsection{10. Measurement 5: the Direction Lottery and
Rectification}\label{measurement-5-the-direction-lottery-and-rectification}

\textbf{The direction lottery} (Figure 9a). The generation direction of
an individual seed (matter side \(+\) or antimatter side \(-\)) is set
by the seed's phase profile. Over 50 random-phase packet seeds,
\(P(+)=0.52\) --- very nearly a 50/50 lottery. The instantaneous sign
predictor built from the odd projection of the vertex rate scores 0.70
--- above chance but short of the preregistered 0.9. This shortfall is
not evidence of randomness; it indicates that the present predictor's
observables likely fail to contain all the state information that
determines the sign --- candidates include hair-resolved phase
alignment, interference among channels, and the projection sign onto the
unstable eigenmode. Constructing a complete instantaneous predictor is
an open problem (Section 14-2). In the channel-multiplicity test, the
generation rate grows with the number of seed harmonics \(n_h\) with
power \(q=0.91\), consistent with the reading that register locality
(Section 9) is the Fourier dual of four-wave-mixing channel count.

\textbf{Ratchet rectification} (Figure 9b). If the direction is a 50/50
lottery, why does net matter remain? The autocatalysis \(Cf^2\) is a
rectifier --- the side that grows grows faster: over a random-phase
ensemble of 20 runs \(\times\) 800 steps, \textbf{80\% grew net} (mean
final ratio 1.0013). Net content arises with no hand-placed symmetry
breaking, from reversible elementary processes plus autocatalysis alone.
This is a measurement of a mechanism inside the model; we do not claim
it explains the matter excess of the real universe.

\begin{figure}
\centering
\pandocbounded{\includegraphics[keepaspectratio,alt={Figure 9 lottery and ratchet}]{figs/fig_mb6_lottery_ratchet_v1.png}}
\caption{Figure 9 lottery and ratchet}
\end{figure}

\textbf{Figure 9}: (a) signed generation rates (six random-phase
examples) and the 50-seed statistics. (b) ensemble trajectories and
their mean (red).

\subsection{11. Recording a Null Result --- E8 Special-Value
Selectivity}\label{recording-a-null-result-e8-special-value-selectivity}

We tested whether register counts \(N_{reg}\) at E8-related special
values (124, 137, 144, 248) behave selectively, using four instrument
families --- all-band linear phases, first-order ignition rates,
thermalization entropy, and ladder-front collision interference
(Nyquist-band power). \textbf{All null} (every special-to-neighbor ratio
inside the preregistered band {[}0.5, 2{]}). We record the motivation
(an internal E8 structural hypothesis of the series), the instrument
designs, and the failed early instruments (Section 12). Establishing
that selectivity is absent deserves the same weight as claiming it is
present.

\subsection{12. Record of Refutations and
Corrections}\label{record-of-refutations-and-corrections}

{\def\LTcaptype{none} % do not increment counter
\begin{longtable}[]{@{}
  >{\raggedright\arraybackslash}p{(\linewidth - 6\tabcolsep) * \real{0.1429}}
  >{\raggedright\arraybackslash}p{(\linewidth - 6\tabcolsep) * \real{0.2857}}
  >{\raggedright\arraybackslash}p{(\linewidth - 6\tabcolsep) * \real{0.2857}}
  >{\raggedright\arraybackslash}p{(\linewidth - 6\tabcolsep) * \real{0.2857}}@{}}
\toprule\noalign{}
\begin{minipage}[b]{\linewidth}\raggedright
\#
\end{minipage} & \begin{minipage}[b]{\linewidth}\raggedright
Error
\end{minipage} & \begin{minipage}[b]{\linewidth}\raggedright
Discovery
\end{minipage} & \begin{minipage}[b]{\linewidth}\raggedright
Correction
\end{minipage} \\
\midrule\noalign{}
\endhead
\bottomrule\noalign{}
\endlastfoot
1 & Conservation criteria measured non-invariants (pointwise closure is
not invariant under slice-independent linear dynamics) with an
unjustified threshold \(10^{-6}\) & U2/U3 FAIL & Derive the invariant
first, then judge (all pass in v2). Discipline: \textbf{always derive
the invariant before preregistering a conservation criterion} \\
2 & The same error recurred (slice closure is not invariant under vertex
ON) & Stage-3 B2 FAIL & Re-judged on the pointwise field closure (the
true invariant); discipline written out \\
3 & Generation-rate predictors wrong twice (cumulative
vs.~instantaneous; own-edge product vs.~mediated product) & D1-P2 FAIL &
The two-channel complete predictor (parametric + beat-XPM) --- stepwise
coherence 0.9998 \\
4 & Blind spots of the D2 instrument (ignition window closed to low
slices; the negative-frequency ladder alias misread as wall arrival) &
\(t_{wrap}=1\) for all \(N_{reg}\) & Third-generation ladder-front
instrument (positive-branch restriction, front tracking) \(\to\) null
confirmed \\
5 & Identically-zero defect of the motion gauge: for odd-harmonic
packets (\(|W|^2\) has period \(N/2\)) the fundamental winding moment
vanishes identically --- the v1 position readings were noise & Found in
figure inspection (\(|z_1|\sim10^{-16}\) confirmed) & Corrected to the
winding-2 moment (position defined mod \(N/2\)); additionally calibrated
a sign error in the prediction (\(e^{+ik\omega_1}\) translates in the
\(-\) direction). After correction: 1D all 8 points within
\(3\times10^{-4}\); 3D all points 0.0000 \\
\end{longtable}
}

The three families (\#1--2, \#3, \#5) are all the same mistake: building
an instrument without deriving the quantity it measures. Having stepped
on the same rake three times, we write the discipline out:
\textbf{derive the invariance (or the predicted value) of the measured
quantity before designing the instrument.}

\subsection{13. Prior Work --- Correspondences Found After Independent
Construction}\label{prior-work-correspondences-found-after-independent-construction}

The architecture of this paper was constructed independently from the
two axioms and the two-body vertex, not designed by relying on the
following. Re-surveying after the design, we found these
correspondences. They are supporting evidence that the framework is not
an isolated fantasy; they are not proofs of correctness.

\begin{enumerate}
\def\labelenumi{\arabic{enumi}.}
\tightlist
\item
  \textbf{Variables on edges --- lattice gauge theory} {[}8{]}. In
  Wilson's lattice gauge theory the gauge field (link variables) lives
  on edges, with interactions built through sites and plaquettes.
  Differences: a background lattice (fixed geometry) and a gauge group
  are inputs there, and matter fields live on sites; here there is no
  geometry (complete graph), groups are post-hoc classifications, and
  nodes carry no variables at all.
\item
  \textbf{Spacetime emergent from pairwise degrees of freedom --- matrix
  models (IKKT/BFSS)} {[}9{]}. The fundamental degrees of freedom are
  matrices whose off-diagonal entries are strings connecting brane pairs
  (pairwise relational degrees of freedom); there is no ambient
  spacetime, and in Lorentzian IKKT Monte Carlo the nine-dimensional
  symmetry breaks spontaneously to an expanding three-dimensional space
  {[}10{]} --- phenomenologically parallel to the three-direction
  emergence of this series. Differences: the action comes from
  superstring theory (supersymmetry as input); matrices carry diagonals
  (per-body position variables) whereas we are purely relational; and
  the address of locality differs (next item).
\item
  \textbf{Complete graphs and locality --- quantum graphity} {[}11{]}.
  Identical to us up to the starting point --- degrees of freedom on a
  complete graph with physics invariant under vertex permutations (=
  anonymity) --- graphity proposed that locality emerges by the graph
  thinning out in a low-energy phase. Our measurement is \textbf{the
  opposite answer to the same question}: the graph stays complete,
  locality does not appear on it (Section 9), and it lives on the
  register axis.
\item
  \textbf{Edge oscillators with node mediation --- topological Kuramoto}
  {[}12{]}. The recent framework placing phase oscillators on edges
  coupled through Hodge Laplacians (node sharing) is the closest
  mathematical form to our adjacency structure. Differences: their
  coupling is gradient-type and the subject is synchronization; ours is
  the cubic exchange flow selected by closure conservation, and the
  subject is generation.
\item
  \textbf{Exact reduction of all-to-all coupling --- Kuramoto mean
  field} {[}13{]}. The order-parameter reduction \(O(N^2)\to O(N)\) is
  the classical precedent of our node ledgers. Differences: ours
  aggregates edge variables at nodes with self-subtraction (four ledgers
  including the anomalous \(z^2\)), and the basis of exactness is
  anonymity (absence of pair-specific constants).
\item
  \textbf{End-to-end numerical genesis --- preheating lattice codes}
  {[}14{]}. LATTICEEASY and its kin follow particle production from
  vacuum in one run. Differences: they input an expanding background, a
  potential, and field content. We have a single anonymous dynamics
  only; spacetime (the three-direction structure) is itself part of the
  output, and so is ``when.''
\item
  \textbf{The relational-dynamics lineage --- Barbour--Bertotti}
  {[}15{]}. The N-body implementation of Mach's principle (dynamics on
  relative configuration space). Difference: there, particle positions
  remain the variables and relationalization is a quotient operation;
  here relations are the variables from the start and body variables
  appear nowhere in the theory.
\end{enumerate}

\subsection{14. Outlook --- Open Problems (working guide; not claimed
here)}\label{outlook-open-problems-working-guide-not-claimed-here}

\begin{enumerate}
\def\labelenumi{\arabic{enumi}.}
\tightlist
\item
  \textbf{Order-parameterization of the condensation gate}: an analytic
  expression for what the gate ratio 2.74 measures. The condensed and
  non-condensed states are matched in norm, seed, and spectral
  occupation, so the difference must lie in phase correlations; the
  candidate is an adjacent-edge relative-phase coherence of the type
  \(\langle \mathrm{Im}(\bar z_e z_{e'})^2\rangle_{adj}\). If extracted,
  the generation law closes as
  \(\Gamma = G(\mathcal{O}_{cond})\cdot Cf^2\).
\item
  \textbf{Projection decomposition of seed unification}: what is shown
  so far is that one \(\delta\) ignited both. Whether
  \(\delta = \delta_\parallel\) (projection onto the unstable
  eigenspace) \(+\ \delta_\perp\) (projection onto the odd-parity
  generation channel) is one degree of freedom or two independent
  components is to be decided by the four-way perturbation experiment
  (eigen-direction on/off \(\times\) odd seed on/off). The complete
  instantaneous sign predictor (Section 10; 0.70 \(\to\) 1) belongs to
  the same projection decomposition.
\item
  \textbf{The thermodynamic limit of the ratchet}: decide whether net
  matter is (a) finite-size fluctuation (mean asymmetry \(\to\) 0 as
  \(N\to\infty\)), (b) spontaneous symmetry selection (final-value
  distribution bimodal at \(\pm m^*\)), or (c) true dynamical
  rectification, by measuring distributions vs.~system size, time, and
  seed number. The present mean final ratio 1.0013 is small; we do not
  speak of matter dominance without this decision.
\item
  \textbf{Marrying the frame and the register}: the unified geometry of
  the three-direction metastable structure (the frame on the graph side)
  and register position; a metric via the chord formula
  \(R''=2R\sin(\Delta\theta/2)\).
\item
  \textbf{Separating parity and hair onto internal axes}: toward pure
  position axes plus internal degrees of freedom (the periodic-table
  program).
\item
  \textbf{E8}: with all instruments null, move the search to the 8-axis
  spinor structure (a separate theoretical line). Finite-order structure
  likely lives not in lattice sizes but in operators, internal axes, and
  spinor closures.
\item
  \textbf{Closed form and antisymmetric structure of the many-body
  flow}: the two-body vertex had a closed-form exact solution (pointwise
  rotation). In the mediated adjacency sum, multiple rotation tangents
  act simultaneously and the individual driving scalars
  \(s_{ee'}=\mathrm{Im}(\bar z_{e'}z_e)\) can change under the other
  adjacent flows, so the closed form is not trivially inherited. Can the
  whole many-body vertex be written as a state-dependent antisymmetric
  flow \(\dot{\mathbf z}=A(\mathbf z)\mathbf z\) with
  \(A^{\mathsf T}=-A\)? If so, the uniqueness theorem advances to the
  geometric form ``closure conservation selects the antisymmetric
  tangent flow on the many-body edge space.'' This is the next
  theoretical stronghold.
\end{enumerate}

\subsection{15. What This Paper Does Not Claim;
Falsifiability}\label{what-this-paper-does-not-claim-falsifiability}

This paper does not claim that the real universe is this model. The
two-factor law, seed unification, and register locality are consequences
of ``if this dynamics and this embedding are assumed.'' Falsification
conditions: (a) generation that bypasses the condensation gate (equal
rates from non-condensed states); (b) non-unified seeds (a divergence of
inflation seed and matter seed violating the \(\varepsilon\)-law); (c)
the discovery of graph locality (concentration systematically changing
generation rates); (d) breakdown of the transplanted laws (\(p\neq2\) or
non-constant \(C\)) --- any of these, demonstrated in a more constrained
experimental system, rejects the embedding claims. The value of the
assumption is measured not by correctness but by the
\textbf{demonstrated connection} --- the physics of two independently
published papers reproducing simultaneously as different chapters of a
single, switchless dynamics.

\textbf{Closing}: the entire content of this paper fits in three lines.
The unstable growth of the linear part saturates and builds a condensed
order (the three-direction metastable structure). That condensed order
opens the generation rate of the always-on nonlinear vertex. Hence ---
how inflation ends is how matter generation begins.

\subsection{16. Reproducibility}\label{reproducibility}

We follow the series convention (self-contained reproduction; originals
imported read-only after snapshot verification; criteria fixed before
execution).

\begin{itemize}
\tightlist
\item
  \textbf{Originals}: Paper 8
  \texttt{run\_preliminary\_seed\_ablation\_v1.py} (SHA-256
  75a10a5b\ldots); generator v3
  \texttt{make\_parent\_white\_harmonics\_n\_only\_v3.py} (SHA-256
  d3217579\ldots).
\item
  \textbf{Code of this paper} (repository folder
  \texttt{universal\_interaction\_manybody\_connection\_v1}, under the
  series folder for dimensional generation structure): the design and
  theorem documents; stage 1--3 engines and tests
  (\texttt{run\_stage1\_multislice\_engine\_v1.py},
  \texttt{run\_stage2\_vertex\_engine\_v1.py},
  \texttt{run\_stage2\_tests\_corrected\_v2.py},
  \texttt{run\_stage3\_sharedO\_v2\_and\_hair\_v1.py}, etc.);
  discriminator experiments
  (\texttt{run\_stage2\_discriminators\_v1.py}); GENESIS
  (\texttt{run\_genesis\_v1.py}, \texttt{\_v2\_local},
  \texttt{\_v3\_register\_local}); the remaining-holes experiments
  (\texttt{run\_holes\_h1\_h2\_v1.py},
  \texttt{run\_h3\_ladder\_front\_v1.py}); the ratchet
  (\texttt{run\_ratchet\_test\_v1.py}); motion
  (\texttt{run\_kinetic\_dispersion\_demo\_v1-v3.py},
  \texttt{run\_genesis\_3d\_demo\_v1-v2.py}); figure generation
  (\texttt{make\_paper\_figures\_multibody\_v1.py} --- experiments whose
  series were not stored are re-run with identical seeds and checked
  against the stored summary values before plotting).
\item
  \textbf{Result JSONs}: every number in the text agrees with the
  corresponding \texttt{*\_result\_*.json}.
\item
  \textbf{Machine verification of the mediated theorem}:
  \texttt{test\_mediated\_vertex\_closure\_v1.py} (T1--T5).
\end{itemize}

\subsection{References}\label{references}

\textbf{Within the series (dependence indicated)}

{[}1{]} N. Kihara, \emph{Generation of Fermionic Structure Is Induced,
Autocatalytic, and Pair-Correlated} (2026), Zenodo. Concept DOI:
10.5281/zenodo.21808091 (the embedded vertex and all transplanted laws;
throughout)

{[}2{]} N. Kihara, \emph{Geometric Rapid Expansion Is Specific to
Unstable Self-Consistent Closures} (2026), Zenodo. Concept DOI:
10.5281/zenodo.21798854 (onset modes, three-direction condensation, the
\(\varepsilon\)-law, the generator; Sections 3, 6, 8)

{[}3{]} N. Kihara, \emph{Paper 7: The Temporal Structure of the N-Body
Relational Closed Wave System} (2026), Zenodo. Concept DOI:
10.5281/zenodo.21543070 (the canonical phenomena of the linear dynamics;
Section 4)

{[}4{]} N. Kihara, \emph{Paper 8: Causal Separation by Two-Stage Seed
Removal} (2026), Zenodo. Concept DOI: 10.5281/zenodo.21614402 (the
direct original of the linear-dynamics code; Sections 4, 16)

{[}5{]} N. Kihara, \emph{Fundamental Axioms of the Anonymous
Equal-Amplitude Composite Wave Model} (2026), Zenodo. Concept DOI:
10.5281/zenodo.21315735 (the two axioms; Section 5)

{[}6{]} N. Kihara, \emph{Future-Phase Position-Acceleration Map in an AB
Two-Body Closed Phase System and the Inverse-Square Law via Harmonic
Closure} (2026), Zenodo. DOI: 10.5281/zenodo.21441081 (the
harmonic-closure dispersion \(\omega_k=k\,\omega_1\); Section 9 only)

{[}7{]} N. Kihara, \emph{The Generative Structure of Fermions} (2026),
Zenodo. Concept DOI: 10.5281/zenodo.21766706 (the original definitions
of the readout \(\theta\), \(R\), and the fermionic mask; Section 4)

\textbf{External}

{[}8{]} K. G. Wilson, ``Confinement of quarks,'' Phys. Rev.~D
\textbf{10}, 2445 (1974)

{[}9{]} N. Ishibashi, H. Kawai, Y. Kitazawa, A. Tsuchiya, ``A large-N
reduced model as superstring,'' Nucl. Phys. B \textbf{498}, 467 (1997);
T. Banks, W. Fischler, S. H. Shenker, L. Susskind, ``M theory as a
matrix model,'' Phys. Rev.~D \textbf{55}, 5112 (1997)

{[}10{]} S.-W. Kim, J. Nishimura, A. Tsuchiya, ``Expanding
(3+1)-dimensional universe from a Lorentzian matrix model,'' Phys.
Rev.~Lett. \textbf{108}, 011601 (2012)

{[}11{]} T. Konopka, F. Markopoulou, L. Smolin, ``Quantum graphity,''
arXiv:hep-th/0611197 (2006); T. Konopka, F. Markopoulou, S. Severini,
``Quantum graphity: a model of emergent locality,'' Phys. Rev.~D
\textbf{77}, 104029 (2008)

{[}12{]} A. P. Millán, J. J. Torres, G. Bianconi, ``Explosive
higher-order Kuramoto dynamics on simplicial complexes,'' Phys.
Rev.~Lett. \textbf{124}, 218301 (2020); A. Arnaudon, R. L. Peach, G.
Petri, P. Expert, ``Connecting Hodge and Sakaguchi-Kuramoto through a
mathematical framework for coupled oscillators on simplicial
complexes,'' Commun. Phys. \textbf{5}, 211 (2022)

{[}13{]} Y. Kuramoto, \emph{Chemical Oscillations, Waves, and
Turbulence} (Springer, 1984)

{[}14{]} G. Felder, I. Tkachev, ``LATTICEEASY: A program for lattice
simulations of scalar fields in an expanding universe,'' Comput. Phys.
Commun. \textbf{178}, 929 (2008); A. V. Frolov, ``DEFROST: A new code
for simulating preheating after inflation,'' JCAP \textbf{0811}, 009
(2008)

{[}15{]} J. B. Barbour, B. Bertotti, ``Mach's principle and the
structure of dynamical theories,'' Proc. R. Soc. A \textbf{382}, 295
(1982)

\end{document}
